\sectionkicker{Source-backed technical notes}
\subsection*{Autonomyの後ろからControl Layerが追う — Safety・Scheming・Prompt Injection: Technical Notes}
本欄は記事本文で圧縮した一次資料上の情報を、比較・再検証しやすい形で整理したものである。時系列、確認済みの主張、留意点、一次資料URLを掲載する。
\medskip
\subsection*{Theme at a glance}
\addcontentsline{toc}{subsection}{Theme at a glance}
\begin{center}
\footnotesize
\begin{tabularx}{\linewidth}{@{}p{0.20\linewidth}p{0.10\linewidth}p{0.14\linewidth}X@{}}
\toprule
資料 & 位置づけ & 種別 & 時系列 \\
\midrule
OpenAI–Anthropic joint safety evaluation & 主要資料 & 安全性関連 & 2025-08-27 (CROSS LAB SAFETY EVALUATION（技術イベント）)  \\
ChatGPT Atlas / OWL / prompt-injection hardening & 補足資料 & Agent & 2025-10-21 (ATLAS（公開）); 2025-10-30 (OWL ARCHITECTURE（公開）); 2025-12-22 (ATLAS SECURITY HARDENING（技術イベント）)  \\
gpt-oss-safeguard & 補足資料 & 安全性関連 & 2025-10-29 (オープンウェイト SAFETY モデル公開（技術イベント）)  \\
Scheming detection and mitigation research & 補足資料 & 安全性関連 & 2025-09-17 (SAFETY 研究（公開）)  \\
\bottomrule
\end{tabularx}
\end{center}
\normalsize

\begin{technicalnote}{OpenAI–Anthropic joint safety evaluation}{主要資料}
\begingroup
% reader-facing Technical Notes break policy
\widowpenalty=10000
\clubpenalty=10000
\displaywidowpenalty=10000
\begin{tabularx}{\linewidth}{@{}>{\bfseries}p{0.22\linewidth}X@{}}
組織 & OpenAI / Anthropic \\
種別 & 安全性関連
時系列 & 2025-08-27 (CROSS LAB SAFETY EVALUATION（技術イベント）) \\
\end{tabularx}
\smallskip
{\bfseries 一次資料から整理したtechnical points}
\begin{itemize}[leftmargin=1.5em,itemsep=0.35em]
\item \textbf{一次情報で確認できる事実}: OpenAIとAnthropicは共同評価を公開し、両社が互いのモデルを対象に、ミスアラインメント、指示追従、ハルシネーション、ジェイルブレイクなどの安全性観点を評価した。
\end{itemize}
\begin{minipage}{\linewidth}
% reader-facing Technical Notes generic boundary/source tail group
{\bfseries 読む際の境界}
\begin{itemize}[leftmargin=1.5em,itemsep=0.35em]
\item \textbf{分析上の留意点}: ここで確認できるのは一次資料に記録された範囲の事実であり、ベンダー・プロジェクト・著者による評価を独立再現された結果として扱うものではない。
\end{itemize}
\begin{samepage}
{\bfseries 一次資料}
\begingroup\sloppy
\Urlmuskip=0mu plus 2mu\relax
\begin{itemize}[leftmargin=1.5em,itemsep=0.25em]
\item {\scriptsize\url{https://openai.com/index/openai-anthropic-safety-evaluation}}
\end{itemize}
\endgroup
\end{samepage}
\end{minipage}
\endgroup
\end{technicalnote}
\medskip

\begin{technicalnote}{ChatGPT Atlas / OWL / prompt-injection hardening}{補足資料}
\begingroup
% reader-facing Technical Notes break policy
\widowpenalty=10000
\clubpenalty=10000
\displaywidowpenalty=10000
\begin{tabularx}{\linewidth}{@{}>{\bfseries}p{0.22\linewidth}X@{}}
組織 & OpenAI \\
種別 & Agent
時系列 & 2025-10-21 (ATLAS（公開）); 2025-10-30 (OWL ARCHITECTURE（公開）); 2025-12-22 (ATLAS SECURITY HARDENING（技術イベント）) \\
\end{tabularx}
\smallskip
{\bfseries 一次資料から整理したtechnical points}
\begin{itemize}[leftmargin=1.5em,itemsep=0.35em]
\item \textbf{一次情報で確認できる事実}: 「ChatGPT Atlas / OWL / prompt-injection hardening」について、一次資料で確認できる公開・提供・機能・時系列上の事実を要約した項目。
\item \textbf{一次情報で確認できる事実}: 「ChatGPT Atlas / OWL / prompt-injection hardening」について、一次資料で確認できる公開・提供・機能・時系列上の事実を要約した項目。
\end{itemize}
\begin{minipage}{\linewidth}
% reader-facing Technical Notes generic boundary/source tail group
{\bfseries 読む際の境界}
\begin{itemize}[leftmargin=1.5em,itemsep=0.35em]
\item \textbf{分析上の留意点}: 「ChatGPT Atlas / OWL / prompt-injection hardening」の扱いでは、一次資料で確認できる範囲、提供時点、評価条件を維持し、記載範囲を超えて一般化しない。
\end{itemize}
\begin{samepage}
{\bfseries 一次資料}
\begingroup\sloppy
\Urlmuskip=0mu plus 2mu\relax
\begin{itemize}[leftmargin=1.5em,itemsep=0.25em]
\item {\scriptsize\url{https://openai.com/index/building-chatgpt-atlas}}
\item {\scriptsize\url{https://openai.com/index/hardening-atlas-against-prompt-injection}}
\item {\scriptsize\url{https://openai.com/index/introducing-chatgpt-atlas}}
\end{itemize}
\endgroup
\end{samepage}
\end{minipage}
\endgroup
\end{technicalnote}
\medskip

\begin{technicalnote}{gpt-oss-safeguard}{補足資料}
\begingroup
% reader-facing Technical Notes break policy
\widowpenalty=10000
\clubpenalty=10000
\displaywidowpenalty=10000
\begin{tabularx}{\linewidth}{@{}>{\bfseries}p{0.22\linewidth}X@{}}
組織 & OpenAI \\
種別 & 安全性関連
時系列 & 2025-10-29 (オープンウェイト SAFETY モデル公開（技術イベント）) \\
\end{tabularx}
\smallskip
{\bfseries 一次資料から整理したtechnical points}
\begin{itemize}[leftmargin=1.5em,itemsep=0.35em]
\item \textbf{一次情報で確認できる事実}: OpenAIはgpt-oss-safeguardを、指定したポリシーに基づく安全性分類を行うオープンウェイト推論モデルとして公開した。
\end{itemize}
\begin{minipage}{\linewidth}
% reader-facing Technical Notes generic boundary/source tail group
{\bfseries 読む際の境界}
\begin{itemize}[leftmargin=1.5em,itemsep=0.35em]
\item \textbf{分析上の留意点}: ここで確認できるのは一次資料に記録された範囲の事実であり、ベンダー・プロジェクト・著者による評価を独立再現された結果として扱うものではない。
\end{itemize}
\begin{samepage}
{\bfseries 一次資料}
\begingroup\sloppy
\Urlmuskip=0mu plus 2mu\relax
\begin{itemize}[leftmargin=1.5em,itemsep=0.25em]
\item {\scriptsize\url{https://openai.com/index/introducing-gpt-oss-safeguard}}
\end{itemize}
\endgroup
\end{samepage}
\end{minipage}
\endgroup
\end{technicalnote}
\medskip

\begin{technicalnote}{Scheming detection and mitigation research}{補足資料}
\begingroup
% reader-facing Technical Notes break policy
\widowpenalty=10000
\clubpenalty=10000
\displaywidowpenalty=10000
\begin{tabularx}{\linewidth}{@{}>{\bfseries}p{0.22\linewidth}X@{}}
組織 & OpenAI / Apollo Research \\
種別 & 安全性関連
時系列 & 2025-09-17 (SAFETY 研究（公開）) \\
\end{tabularx}
\smallskip
{\bfseries 一次資料から整理したtechnical points}
\begin{itemize}[leftmargin=1.5em,itemsep=0.35em]
\item \textbf{一次情報で確認できる事実}: OpenAIとApollo Researchは、隠れたミスアラインメント、いわゆる「scheming」を対象とする管理された評価と、初期的な緩和手法を公開した。
\end{itemize}
\begin{minipage}{\linewidth}
% reader-facing Technical Notes generic boundary/source tail group
{\bfseries 読む際の境界}
\begin{itemize}[leftmargin=1.5em,itemsep=0.35em]
\item \textbf{分析上の留意点}: これらの所見は管理された評価環境で得られたものであり、記載された実験範囲を超えて実運用モデルの挙動へ一般化しない。
\end{itemize}
\begin{samepage}
{\bfseries 一次資料}
\begingroup\sloppy
\Urlmuskip=0mu plus 2mu\relax
\begin{itemize}[leftmargin=1.5em,itemsep=0.25em]
\item {\scriptsize\url{https://openai.com/index/detecting-and-reducing-scheming-in-ai-models}}
\end{itemize}
\endgroup
\end{samepage}
\end{minipage}
\endgroup
\end{technicalnote}
\medskip
